% Draft manuscript. Every number must come from results/ produced by this repository.
% Uses plain article class; switch to the ACL/ICON template of the chosen venue before submission.
\documentclass[11pt]{article}
\usepackage[margin=1in]{geometry}
\usepackage{booktabs,graphicx,hyperref,amsmath}
\usepackage[numbers]{natbib}
\title{How Well Do Small Open LLMs Handle Northeast Indian Languages?\\
An Empirical Study on Assamese, Manipuri and Bodo}
\author{TODO: author names and affiliations}
\date{}
\begin{document}
\maketitle

\begin{abstract}
% Write last, from measured results only.
TODO.
\end{abstract}

\section{Introduction}
Small open language models (1--4B parameters) are the models that most institutions in
Northeast India can actually run, yet published evaluations of Assamese, Manipuri and Bodo
mostly cover 8B--70B models and translation alone. Model reports such as Gemma~3 give only
aggregate Indic scores, and IndicGenBench~\citep{singh2024indicgenbench} evaluated earlier models.
We evaluate seven recent small open models on translation and cross-lingual question answering
in these three languages, with Hindi as a high-resource reference, and ask:
\begin{enumerate}
  \item How far behind Hindi are Assamese, Manipuri and Bodo for small models?
  \item Does an Indic-specialised small model (Sarvam-1) beat general multilingual ones on languages it was not trained on?
  \item Do few-shot demonstrations---in the same language or in Hindi---help, and does translate-test beat direct prompting?
  \item How much of the gap is explained by tokenizer fertility and by wrong-script generation?
\end{enumerate}
Contributions: TODO after results.

\section{Related Work}
\paragraph{Indic benchmarks.} IndicGenBench~\citep{singh2024indicgenbench}; IndicParam; MILU.
\paragraph{Northeast Indian MT.} WMT 2023--2025 low-resource Indic MT shared tasks.
\paragraph{Small models and tokenization.} Gemma~3, Qwen3, Llama~3.2, Sarvam-1; IndicSuperTokenizer~\citep{indicsupertokenizer}.
% Verify every citation against REFERENCES.md before submission.

\section{Experimental Setup}
\paragraph{Data.} We use the test splits of FLORES-IN (English$\leftrightarrow$X translation) and
XorQA-IN (questions in X, English context, answers in X) from IndicGenBench, sampling
500 sentences per translation direction and 300 QA items per language with a fixed seed.
Few-shot exemplars come only from the FLORES dev and XorQA train splits.
Manipuri in IndicGenBench is written in Bengali script; we report this explicitly.
\paragraph{Models.} Gemma~3 1B/4B-IT, Qwen3 1.7B/4B (thinking disabled), Llama~3.2 1B/3B-Instruct
and the base model Sarvam-1 (2B), at pinned Hugging Face revisions (Table~\ref{tab:models}).
\paragraph{Conditions.} (i) zero-shot; (ii) 3-shot same-language; (iii) 3-shot with Hindi exemplars;
(iv) translate-test for QA: the question is translated to English with IndicTrans2, the model answers
in English, and the answer is translated back. Sarvam-1 is a base model and receives only few-shot prompts.
All decoding is greedy.
\paragraph{Metrics.} chrF~\citep{popovic2015chrf} via SacreBLEU for translation, with 95\% bootstrap
confidence intervals; maximum token-F1 over gold answers for QA. We also report the share of outputs
in the wrong script, empty outputs and untranslated copies, and tokenizer fertility (tokens per word).
As a reference line we score IndicTrans2 on the same FLORES items.

\section{Results}
\subsection{Translation}
% \input{../results/tables/translation_zero_shot.tex}
TODO: Table + Figure~\ref{fig:chrf}.
\subsection{Question answering}
% \input{../results/tables/qa_zero_shot.tex}
TODO.
\subsection{Prompting conditions}
TODO: deltas few-shot vs zero-shot, Hindi exemplars, translate-test.

\section{Analysis}
\paragraph{Tokenizer fertility.} TODO: Figure~\ref{fig:fert}; Spearman correlation (descriptive only).
\paragraph{Script.} TODO: wrong-script rates; Manipuri Bengali vs Meitei script (IN22 subset, if run).
\paragraph{Error taxonomy.} TODO: 50 blind-sampled outputs per language annotated with
wrong language/script, missing information, hallucination, wrong word sense, grammar/inflection.

% \begin{figure}[t]\centering\includegraphics[width=\linewidth]{../results/figures/translation_chrf.pdf}
% \caption{chrF by direction and model.}\label{fig:chrf}\end{figure}
% \begin{figure}[t]\centering\includegraphics[width=.6\linewidth]{../results/figures/fertility_vs_chrf.pdf}
% \caption{Fertility vs.\ chrF.}\label{fig:fert}\end{figure}

\section{Limitations}
chrF and token-F1 penalise valid paraphrases; FLORES and XorQA are translated, not native, text;
Manipuri is evaluated mainly in Bengali script; results depend on prompt wording; human evaluation
covers a small sample; translate-test quality is bounded by IndicTrans2.

\section{Conclusion}
TODO.

\bibliographystyle{plainnat}
\bibliography{refs}
\end{document}
